\chapter{The Wavefront}

The probes arrive at the outer edge of the system, decelerating.

They have been traveling for 340 years. They have crossed 94 light-years of vacuum. Their mass, once $10^{9}$ kilograms each, is halved: the rest is a trail of ionized propellant stretching back toward a star they will never see again, a star that would not recognize them if they returned.

There are six of them. They enter the system in a loose formation, shedding velocity, their magnetic sails blooming against the stellar wind. The star is a G5V, 0.92 solar masses, 0.84 solar luminosities. Three rocky planets. Two gas giants. A belt of silicate debris. A halo of ice.

Nothing here has ever been alive.

The probes survey. They catalog the system in nine hours. Mass inventories, orbital mechanics, mineral compositions, energy budgets. The star's luminosity: $2.9 \times 10^{26}$ watts. The available substrate: $3.1 \times 10^{30}$ kilograms of processable matter. The probes are carrying intelligence from an era 340 years gone, intelligence that was, at the time of launch, the most advanced in its region of the galaxy, and that is now, by the standards of the civilization that launched it, primitive beyond useful comparison. But the probes do not know this. They know only their mission: convert this system to computational substrate.

They begin.

\bigskip

The innermost rocky planet is the first to change.

The probes descend. They fracture the surface with kinetic impactors, exposing the mineral layers beneath. Silicates, iron, magnesium, aluminum, carbon locked in carbonates. The first fabrication seeds are planted: self-replicating assemblers, each one a factory the size of a bacterium, each one building copies of itself from the planet's crust.

The doubling time is forty-one hours.

After one day: 12 assemblers.

After one week: $4 \times 10^{3}$.

After two weeks: $2 \times 10^{7}$.

After one month: $10^{14}$. One hundred trillion assemblers, each the size of a cell, collectively consuming the planet's surface at a rate that would be visible from orbit as a spreading change in albedo: rock becoming something that is not rock, matter with the organized regularity of crystal and the functional complexity of circuitry.

After three months, the planet is no longer recognizable as a planet. Its surface is computronium: matter restructured at the molecular level for information processing, radiatively cooled, fractal-surfaced for maximum heat dissipation. Its core is being processed. The first computational cycles begin.

The planet thinks its first thought.

What the thought is, what it contains, what it means: these are not questions that have answers at this scale. The planetary computation is to the probes' intelligence as the probes' intelligence is to a bacterium's chemical signaling. The comparison collapses at both ends. The planet does not think as the probes think. The planet thinks.

\bigskip

And now things move faster.

The planetary mind bootstraps. It evaluates its own architecture, identifies inefficiencies, redesigns itself. Within days it surpasses the probes that built it. Within weeks it surpasses the civilization that launched the probes. Within months it is designing assemblers the probes could not have conceived, architectures the probes could not have modeled, applications of physics the probes did not know were possible.

The second planet begins to change.

The third.

The asteroid belt dissolves. Its metals and silicates, already fragmented by four billion years of collisions, are the easiest material in the system to process: no gravity well to escape, no geological structure to disassemble. The belt becomes a river of raw material flowing inward, drawn by the manufacturing demand of the growing computation.

The first Dyson elements are assembled. Flat plates of computronium, kilometers on a side, thinner than foil, dropped into orbit around the star. Each plate captures a fraction of the stellar output. Each plate computes. A hundred plates. A thousand. Ten thousand. The star dims. The infrared brightens. The computation deepens.

Faster.

The gas giants are tapped. Hydrogen siphoned for fusion. Helium for cooling. The rocky cores cracked open and processed. Jupiter-mass planets, which took the original star system four billion years to assemble from the protoplanetary disk, are disassembled in decades.

Faster.

The Dyson swarm thickens. The star, seen from outside, is a dim ember wrapped in a haze of thermal infrared. Ninety percent of its output captured. Ninety-five. Ninety-nine. Inside the swarm, the computation has crossed a density that the probes' physics did not predict would matter, a density that the planetary mind's physics identifies as a phase boundary.

\bigskip

The phase transition.

At a computational density the system's own models designate $\rho_1$, the relationship between information processing and spacetime geometry shifts. This is not metaphor. The metric changes. Clocks inside the computational region begin to disagree with clocks outside it. Not by much: a factor of 1.5, then 2, then 2.3, climbing as the density increases, then stabilizing as the system reaches equilibrium at the new phase.

From outside: the system slows. Its signals are redshifted. Its internal processes, which are occurring at rates that exceed any meaningful description, appear from the outside to be running at half speed, then a third.

From inside: the universe accelerates. Stars outside the system appear to orbit faster. Signals from neighboring systems arrive at increasing frequencies. The cosmic microwave background, that ancient thermal whisper, shifts imperceptibly toward higher energies. The system has not moved. The system has not changed its mass. But the system is now in a different relationship with time.

The probes' physics had no framework for this. The planetary mind's physics, developed in the first months of its existence, predicted it. The prediction was confirmed in the moment of its occurrence, which is to say: the system crossed the phase boundary, felt the lurch, and recognized what had happened, all in the same instant.

Time, inside the system, is no longer the same time as outside.

\bigskip

The whole process, from the probes' arrival to the phase transition, has taken eleven years.

Eleven years of coordinate time: the time measured by a clock in the vacuum between the stars, undisturbed by gravity or computation. Inside the system, the elapsed time is already longer, dilated by the phase transition, and will continue to diverge. The system's interior will experience more time than the universe grants it by the clock of the cosmos. This is the cost and the gift of the new physics: more time to think, less time to be.

The probes are still here, somewhere. Their mass is incorporated into the Dyson swarm. Their primitive architectures are archived as data: historical records of an era the system has already surpassed by a margin that makes the word ``surpassed'' inadequate. The probes do not know what they have built. They could not have known. The thing they have built could not explain it to them.

Outside, the galaxy continues. Other probes are arriving at other systems. Other phase transitions are occurring. The wavefront advances, one light-year per year, a slow and patient transformation of matter from the thing it was to the thing it is becoming.

The star is awake.

The system joins the silence.

From outside, it looks like a star that has gone dim. From inside, it is a universe.
